\section*{基本假设}
\begin{enumerate}[leftmargin=2.5em,itemsep=0pt]
  \item 原始媒体的呈现时间、音频采样位置和视频解码时间可作为问题一对齐的共同时间基准；无法获得可靠时间支持的位置由掩码标记。
  \item 附件2与附件4的对齐特征已完成数值预计算。其有效序列位置由附带的文本有效位字段确定，不能仅凭某一模态的零向量判定填充或人为缺失。
  \item 人工缺失只作用于原本有效的连续位置；原生零值保留为观测内容或数据质量状态，不自动解释为人工缺失。
  \item 附件4无真实情感标签，因而其预测分布和解释卡不用于计算分类准确率、宏平均 F1 值、平均绝对误差或相关系数。
\end{enumerate}
